\section*{अध्याय \ref{ch:b2:planet-formation-meteorites} --- ग्रहों का निर्माण और उल्कापिंड}

\begin{solution}{exo:b2:planet-formation-meteorites:1}
गैस: \qty{1}{au} पर \qty{2200}{g.cm^{-2}}; \qty{5.2}{au} पर,
$2200\times5.2^{-3/2} = 2200/11.86 = \qty{186}{g.cm^{-2}}$। ठोस:
\qty{1}{au} पर $33\times0.33 = \qty{10.9}{g.cm^{-2}}$;
\qty{5.2}{au} पर $33/11.86 = \qty{2.78}{g.cm^{-2}}$। \qty{5.2}{au} पर गैस ठोसों से
$186/2.78 = 67$ गुना भारी है।
\end{solution}

\begin{solution}{exo:b2:planet-formation-meteorites:2}
ऐलुमिनियम (\qty{1653}{K}, अति-दुर्गलनीय), मैग्नीशियम (\qty{1336}{K}, सामान्य),
लोहा (\qty{1334}{K}, सामान्य), सोडियम (\qty{958}{K}, मध्यम वाष्पशील), सल्फ़र
(\qty{664}{K}, वाष्पशील), जल-बर्फ़ के रूप में ऑक्सीजन (\qty{180}{K}, अति वाष्पशील)।
CAI, जो प्रथम संघनित पदार्थ हैं, ऐलुमिनियम (और कैल्सियम, टाइटेनियम) में समृद्ध हैं;
पृथ्वी मैग्नीशियम और ऐलुमिनियम की तुलना में वाष्पशील तत्त्वों, सोडियम, सल्फ़र और जल,
में न्यून है।
\end{solution}

\begin{solution}{exo:b2:planet-formation-meteorites:3}
\omterm{def:b2:planet-formation-meteorites:fall-find}{पतित उल्कापिंड}: साधारण \omterm{def:b2:planet-formation-meteorites:chondrule}{कॉन्ड्राइट} $856/1107 = \qty{77.3}{\percent}$, लौह
$49/1107 = \qty{4.4}{\percent}$। \omterm{def:b2:planet-formation-meteorites:fall-find}{प्राप्त उल्कापिंड}: $38647/44534 = \qty{86.8}{\percent}$ और
$1020/44534 = \qty{2.3}{\percent}$। अंटार्कटिक \omterm{def:b2:planet-formation-meteorites:fall-find}{प्राप्त उल्कापिंड}: $22096/44534 =
\qty{49.6}{\percent}$। \omterm{def:b2:planet-formation-meteorites:fall-find}{पतित उल्कापिंड}: हर प्रेक्षित पतन एकत्र किया जाता है, उसका
वर्ग जो भी हो, जबकि प्राप्ति इस पर निर्भर करती है कि \omterm{def:b2:planet-formation-meteorites:fall-find}{उल्कापिंड} भूमि पर कितनी अच्छी तरह बचा रहता है,
वह कैसा दिखता है और लोग कहाँ खोजते हैं, और आधे \omterm{def:b2:planet-formation-meteorites:fall-find}{प्राप्त उल्कापिंड} एक ही प्रकार के
स्थान से आते हैं।
\end{solution}

\begin{solution}{exo:b2:planet-formation-meteorites:4}
$\Delta t = 1.035\ln(\num{5.23e-5}/\num{1.0e-5}) = 1.035\ln 5.23 =
1.035\times1.654 = \qty{1.71}{Myr}$, CAI के बाद।
\end{solution}

\begin{solution}{exo:b2:planet-formation-meteorites:5}
$4\pi\Sigma_1 r_1^2 = 4\pi\times\num{2.2e4}\times(\num{1.496e11})^2 =
\qty{6.19e27}{kg}$। \qty{30}{au} के भीतर: $\times\sqrt{30} = 5.48$, अतः
\qty{3.39e28}{kg} $= 0.017$ सौर द्रव्यमान; \qty{100}{au} के भीतर: $\times10$,
\qty{6.19e28}{kg} $= 0.031$ सौर द्रव्यमान। दोनों प्रेक्षित परास के भीतर हैं,
उसके ऊपरी सिरे के निकट: नीहारिका एक साधारण चक्रिका थी।
\end{solution}

\begin{solution}{exo:b2:planet-formation-meteorites:6}
$r = (120/T)^{7/3}$ au: जल के लिए, $(120/182)^{7/3} = 0.659^{7/3} =
\qty{0.38}{au}$; मीथेन के लिए, $(120/41)^{7/3} = 2.93^{7/3} = \qty{12.3}{au}$।
सीधी धूप वाला कण चक्रिका की पृष्ठ-परत है; मध्य-तल,
जो ऊपर की धूल से परिरक्षित है, अधिक ठंडा है, इसलिए बर्फ़ और भीतर तक बची रहती है।
\omterm{def:b2:planet-formation-meteorites:planetesimal}{ग्रहाणु} मध्य-तल में बने, इसलिए मध्य-तल वाला अनुमान लागू होता है, बशर्ते
चक्रिका निष्क्रिय हो; अभिवृद्धि से ऊष्मन \omterm{def:b2:planet-formation-meteorites:snow-line}{हिम रेखा} को बाहर खिसका देता है।
\end{solution}

\begin{solution}{exo:b2:planet-formation-meteorites:7}
\qty{1}{au} पर: $r^2 = \num{2.238e22}$, $\Sigma_p = \qty{109}{kg.m^{-2}}$, अतः
$2\pi b r^2\Sigma_p = 31.42\times\num{2.238e22}\times109 = \num{7.66e25}$;
$(\num{7.66e25})^{3/2} = \num{6.71e38}$, जिसे
$(3\times\num{1.989e30})^{1/2} = \num{2.443e15}$ से भाग देने पर: \qty{2.75e23}{kg}
$= 0.046$ पृथ्वी-द्रव्यमान। \qty{5.2}{au} पर: $r^2 = \num{6.05e23}$,
$\Sigma_p = \qty{27.8}{kg.m^{-2}}$, $2\pi b r^2\Sigma_p = \num{5.29e26}$,
$M = \num{1.22e40}/\num{2.443e15} = \qty{4.98e24}{kg} = 0.83$ पृथ्वी-द्रव्यमान।
\num{3.5} गुना अधिक ठोसों के साथ $M$ $3.5^{3/2} = 6.55$ गुना बढ़ता है: 5.5
पृथ्वी-द्रव्यमान।
\end{solution}

\begin{solution}{exo:b2:planet-formation-meteorites:8}
प्रति किलोग्राम $\num{1.01e19}$ परमाणु, अतः
$\num{1.01e19}\times\num{5.00e-13} = \qty{5.05e6}{J.kg^{-1}}$। तापमान-वृद्धि
$\num{5.05e6}/837 = \qty{6030}{K}$ होती। \num{250} से
\qty{1700}{K} तक गरम करने में $837\times1450 = \qty{1.21e6}{J.kg^{-1}}$ लगते हैं: ऐलुमिनियम-26
ने इसका लगभग चार गुना दिया।
\end{solution}

\begin{solution}{exo:b2:planet-formation-meteorites:9}
$\num{5.07e10}/\num{7.49e7} = 677$। लौह-60 की माध्य आयु
$2.62/\ln2 = \qty{3.78}{Myr}$ है। प्रति ग्राम आरंभिक शक्तियाँ:
$\num{5.07e10}/1.035 = \num{4.90e10}$ और $\num{7.49e7}/3.78 = \num{1.98e7}$ अर्ग प्रति
ग्राम प्रति मिलियन वर्ष, लगभग \num{2470} का अनुपात: ऐलुमिनियम-26 कहीं अधिक प्रबल था।
\end{solution}

\begin{solution}{exo:b2:planet-formation-meteorites:10}
शक्तियाँ तब बराबर होती हैं जब $\num{4.90e10}\,e^{-t/1.035} = \num{1.98e7}\,e^{-t/3.78}$,
अर्थात् $t(1/1.035 - 1/3.78) = \ln(\num{4.90e10}/\num{1.98e7}) = \ln 2475 = 7.81$;
$0.966 - 0.265 = 0.702$ प्रति मिलियन वर्ष के साथ, $t = \qty{11.1}{Myr}$। तब तक
लौह-60 के पास $\num{7.49e7}\,e^{-11.1/3.78} = \num{7.49e7}\times0.053 =
\qty{4.0e6}{erg.g^{-1}}$ बचता है, अर्थात् \qty{400}{J.kg^{-1}}: कुछ भी पिघलाने के लिए हज़ारों
गुना कम।
\end{solution}

\begin{solution}{exo:b2:planet-formation-meteorites:11}
$\tau = 1.035\times\num{3.156e13} = \qty{3.267e13}{s}$ और
$H_0\tau = \qty{6.21e6}{J.kg^{-1}}$। \qty{1200}{K} के लिए:
$t_a^* = 1.035\ln\bigl(\num{6.21e6}/(837\times950)\bigr) = 1.035\ln7.81 =
\qty{2.13}{Myr}$; \qty{1700}{K} के लिए: $1.035\ln\bigl(\num{6.21e6}/(837\times1450)\bigr)
= 1.035\ln5.12 = \qty{1.69}{Myr}$। $H_0$ को आधा करने से दोनों में से $\tau\ln2 =
\qty{0.72}{Myr}$ घट जाता है, क्योंकि $t_a^* = \tau\ln H_0 + \tau\ln(\tau/c(T_c -
T_0))$ और केवल पहला पद बदलता है: खिसकाव एक अर्ध-आयु है, क्रांतिक तापमान
जो भी हो।
\end{solution}

\begin{solution}{exo:b2:planet-formation-meteorites:12}
$\kappa = 2.1/(3300\times837) = \qty{7.60e-7}{m^2.s^{-1}}$ और
$\sqrt{\kappa\tau} = \sqrt{\num{7.60e-7}\times\num{3.267e13}} = \qty{4980}{m}$,
लगभग \qty{5}{km}। \qty{100}{km} चौड़े पिंड का केंद्र \qty{50}{km} गहरा है,
इस दूरी का दस गुना: ऐलुमिनियम-26 के क्षय के दौरान उसकी ऊष्मा बाहर नहीं निकल सकती, और
प्रतिज्ञप्ति लागू होती है। \qty{2}{km} चौड़ा पिंड अपनी ऊष्मा उतनी ही तेज़ी से खोता है जितनी तेज़ी से
बनाता है। गलन के लिए \qty{5}{km} की कुछ गुना त्रिज्या चाहिए, अर्थात् लगभग
\qtyrange{10}{20}{km} या अधिक का व्यास, जैसा मॉडल पाते हैं।
\end{solution}

\begin{solution}{pb:b2:planet-formation-meteorites:1}
\textbf{1.} \omterm{def:b2:planet-formation-meteorites:fall-find}{पतित उल्कापिंड}: उसका आगमन देखा गया और अगली सुबह उसे उठा लिया गया।
वह अपक्षयित नहीं है, और पृथ्वी तक पहुँचने वाली सामग्री के निष्पक्ष नमूने में जुड़ता है।

\textbf{2.} $856/1107 = 0.77$: लगभग चार में से तीन संभावनाएँ।

\textbf{3.} \qty{0.3}{mm} चौड़े \omterm{def:b2:planet-formation-meteorites:chondrule}{कॉन्ड्रयूल}, \qty{9}{\percent} धातु और
\qty{18}{\percent} फ़ेयालाइट वाला ओलिवीन: एक H \omterm{def:b2:planet-formation-meteorites:chondrule}{कॉन्ड्राइट} (\qty{10}{\percent} धातु, Fa 16--20)।

\textbf{4.} धुँधले \omterm{def:b2:planet-formation-meteorites:chondrule}{कॉन्ड्रयूल} और एकसमान खनिज तापीय कायांतरण दर्शाते हैं: प्रकार
5 से 6। साधारण \omterm{def:b2:planet-formation-meteorites:chondrule}{कॉन्ड्राइट} अधिकतम लगभग \qty{1220}{K} तक पहुँचे, अर्थात्
\qty{947}{\degreeCelsius}।

\textbf{5.} साधारण \omterm{def:b2:planet-formation-meteorites:chondrule}{कॉन्ड्राइटों} में आयतन के अनुसार केवल \qtyrange{0.1}{1}{\percent} CAI
होते हैं। पत्थर CAI के बाद एकत्र हुए \omterm{def:b2:planet-formation-meteorites:chondrule}{कॉन्ड्रयूलों} और आधात्री से बना है, और CAI
सबसे पुराने ठोस हैं: उन्हें समाविष्ट करने वाली कोई भी चट्टान उनसे युवा है।

\textbf{6.} हर खनिज को निर्माण के समय वही मैग्नीशियम मिला, साथ में अपने ऐलुमिनियम
के अनुपात में ऐलुमिनियम-26; ऐलुमिनियम-26 तब से
मैग्नीशियम-26 में बदल चुका है, इसलिए किसी खनिज में प्रति मैग्नीशियम जितना अधिक ऐलुमिनियम था,
उसे उतना अधिक मैग्नीशियम-26 मिला।

\textbf{7.} $(0.139590 - 0.139500)/(15.0 - 0.05) = \num{0.000090}/14.95 =
\num{6.02e-6}$: बनते समय \omterm{def:b2:planet-formation-meteorites:chondrule}{कॉन्ड्रयूल} का अनुपात
$\ce{^{26}Al}/\ce{^{27}Al}$।

\textbf{8.} $0.139500 + \num{6.02e-6}\times3.00 = 0.139518$: रेखा पर।

\textbf{9.} $\Delta t = 1.035\ln(\num{5.23e-5}/\num{6.02e-6}) = 1.035\ln8.69 =
1.035\times2.162 = \qty{2.24}{Myr}$।

\textbf{10.} $4567.5 - 2.2 = \qty{4565.3}{\Ma}$।

\textbf{11.} $\num{5.23e-5}\,e^{-10/1.035} = \num{5.23e-5}\times\num{6.4e-5} =
\num{3.3e-9}$, मापने के लिए मैग्नीशियम-26 बहुत ही कम। हैफ़नियम-182 (अर्ध-आयु
\qty{8.9}{Myr}) या मैंगनीज़-53 (\qty{3.7}{Myr}) CAI के लगभग दस मिलियन वर्ष बाद की घटनाओं के
लिए उपयुक्त हैं।

\textbf{12.} $\kappa = 2.1/(3300\times837) = \qty{7.60e-7}{m^2.s^{-1}}$;
$R^2/\kappa = \num{4e8}/\num{7.60e-7} = \qty{5.26e14}{s} = \qty{16.7}{Myr}$, ऐलुमिनियम-26 की
माध्य आयु का सोलह गुना: केंद्र अपनी ऊष्मा बनाए रखता है।

\textbf{13.} $\num{1.9e-7}\,e^{-2.4/1.035} = \num{1.9e-7}\times0.0984 =
\qty{1.87e-8}{W.kg^{-1}}$।

\textbf{14.} $\num{1.87e-8}\times\num{3.267e13} = \qty{6.11e5}{J.kg^{-1}}$।

\textbf{15.} $\Delta T = \num{6.11e5}/837 = \qty{730}{K}$; शिखर
$250 + 730 = \qty{980}{K}$।

\textbf{16.} \qty{750}{K} से ऊपर, इसलिए चट्टान सिंटरित और कायांतरित हुई;
\qty{1200}{K} से नीचे, इसलिए न धातु पिघली न सल्फ़ाइड। पत्थर एक गरम हुआ, उच्च \omterm{def:b2:planet-formation-meteorites:type}{शैलिकी प्रकार} का
\omterm{def:b2:planet-formation-meteorites:chondrule}{कॉन्ड्राइट} है, प्रश्न 4 के \qty{1220}{K} से नीचे: संगत।

\textbf{17.} $1.035\ln\bigl(\num{6.21e6}/(837\times950)\bigr) = \qty{2.13}{Myr}$:
लगभग \qty{0.3}{Myr} पहले।

\textbf{18.} वृद्धि $\Delta T(1 - e^{-(t - t_a)/\tau})$ है; \qty{90}{\percent}
के लिए $e^{-(t - t_a)/\tau} = 0.1$ चाहिए, $t - t_a = 1.035\ln10 = \qty{2.38}{Myr}$, अर्थात्
CAI के \qty{4.8}{Myr} बाद।

\textbf{19.} $\Delta T = 7420\,e^{-0.8/1.035} = 7420\times0.462 = \qty{3430}{K}$,
शिखर \qty{3680}{K}, उस \qty{1700}{K} से बहुत ऊपर जिस पर सिलिकेट गलित होते हैं।
पिंड पिघला (गलन की ऊष्मा ने उसके तापमान को नीचे रखा) और
एक धात्विक क्रोड तथा एक सिलिकेट प्रावार और भूपर्पटी में विभेदित हुआ।

\textbf{20.} $e^{(2.4 - 0.8)/1.035} = e^{1.546} = 4.7$ गुना अधिक।

\textbf{21.} जस्ता और सीसा \qty{727}{K} के निकट संघनित होते हैं, मध्यम वाष्पशील; मैग्नीशियम
(\qty{1336}{K}) और ऐलुमिनियम (\qty{1653}{K}) बहुत पहले संघनित होते हैं। पत्थर के क्षेत्र के
ठोस, गैस के इतना ठंडा होने से पहले ही एकत्र हो गए कि सारा जस्ता और
सीसा संघनित हो सके।

\textbf{22.} $33\times0.33\times2.5^{-3/2} = 10.89/3.95 = \qty{2.76}{g.cm^{-2}}$,
\qty{27.6}{kg.m^{-2}}।

\textbf{23.} $r^2 = (3.74\times10^{11})^2 = \num{1.40e23}$;
$2\pi\times5\times\num{1.40e23}\times27.6 = \num{1.21e26}$;
$(\num{1.21e26})^{3/2}/\num{2.443e15} = \qty{5.45e23}{kg}$, $0.091$ पृथ्वी-द्रव्यमान।
जनक पिंड, $\tfrac43\pi(\num{2e4})^3\times3300 = \qty{1.1e17}{kg}$, पचास
लाख गुना हल्का है: एक \omterm{def:b2:planet-formation-meteorites:planetesimal}{ग्रहाणु}, भ्रूण नहीं।

\textbf{24.} जनक पिंड के केंद्र का शिखर तापमान लगभग
\textbf{\qty[text-series-to-math=true]{707}{\degreeCelsius}} ($980 - 273$) था: \omterm{def:b2:planet-formation-meteorites:chondrule}{कॉन्ड्राइट} को
पुनःक्रिस्टलित करने के लिए पर्याप्त गरम, उसकी धातु को पिघलाने के लिए बहुत ठंडा।
\end{solution}
